\clearpage
\sectionkicker{Source-backed technical notes}
\subsection*{Agents — 長時間動作するシステムへ: Technical Notes}
この欄は記事本文で圧縮した一次資料上の情報を、比較・再検証しやすい形へ展開したものである。新しい外部情報は追加せず、Selection済みEvidenceのchronology、normalized claim、limitations、source URLのみを再配置する。
\medskip
\subsection*{Theme at a glance}
\addcontentsline{toc}{subsection}{Theme at a glance}
\begin{center}
\footnotesize
\begin{tabularx}{\linewidth}{@{}p{0.20\linewidth}p{0.10\linewidth}p{0.14\linewidth}X@{}}
\toprule
Artifact & Role & Type & Objective chronology \\
\midrule
July long-horizon agent evaluation research & PRIMARY & BENCHMARK & 2026-07-31 (PAPER\_PUBLICATION); 2026-07-31 (PAPER\_PUBLICATION) \\
July agent-serving systems research & SUPPORTING & OTHER & 2026-07-13 (PAPER\_PUBLICATION); 2026-07-27 (PAPER\_PUBLICATION); 2026-07-31 (PAPER\_PUBLICATION) \\
\bottomrule
\end{tabularx}
\end{center}
\normalsize

\begin{technicalnote}{July long-horizon agent evaluation research}{PRIMARY}
\begin{tabularx}{\linewidth}{@{}>{\bfseries}p{0.22\linewidth}X@{}}
Organization & - \\
Artifact type & BENCHMARK \\
Chronology & 2026-07-31 (PAPER\_PUBLICATION); 2026-07-31 (PAPER\_PUBLICATION) \\
\end{tabularx}
\smallskip
{\bfseries 一次資料から整理したtechnical points}
\begin{itemize}[leftmargin=1.5em,itemsep=0.35em]
\item \textbf{Author claim}: LoopsBench defines a long-horizon coding-agent benchmark with 112 dependency-structured tasks spanning eight programming languages and nine domains.
\item \textbf{Author claim}: AgentStream evaluates self-evolving agents under isolated, sequential, and interleaved streaming-task settings rather than only independent tasks.
\end{itemize}
{\bfseries 読む際の境界}
\begin{itemize}[leftmargin=1.5em,itemsep=0.35em]
\item \textbf{分析上の留意点}: Both benchmark designs and reported model rankings/results are author-reported; conclusions depend on the included harnesses, models, task sources, and scoring protocols.
\end{itemize}
{\bfseries Primary source}
\begin{itemize}[leftmargin=1.5em,itemsep=0.25em]
\item \url{http://arxiv.org/abs/2608.00155v1}
\item \url{http://arxiv.org/abs/2608.00267v1}
\end{itemize}
{\scriptsize\color{SurveyMuted}Source-bound record: \texttt{evidence:SP-2026-M07:series-agent-eval-july-ebc09433f5}.}
\end{technicalnote}
\medskip

\begin{technicalnote}{July agent-serving systems research}{SUPPORTING}
\begin{tabularx}{\linewidth}{@{}>{\bfseries}p{0.22\linewidth}X@{}}
Organization & - \\
Artifact type & OTHER \\
Chronology & 2026-07-13 (PAPER\_PUBLICATION); 2026-07-27 (PAPER\_PUBLICATION); 2026-07-31 (PAPER\_PUBLICATION) \\
\end{tabularx}
\smallskip
{\bfseries 一次資料から整理したtechnical points}
\begin{itemize}[leftmargin=1.5em,itemsep=0.35em]
\item \textbf{Author claim}: AAFLOW+ treats reusable KV cache as a distributed workflow object for multi-agent pipelines rather than a request-local optimization.
\item \textbf{Author claim}: SpecBox proposes speculative sandbox preallocation/prewarming to trade resource use against cold-start tail latency in sandboxed agent workloads.
\item \textbf{Author claim}: TokTier proposes exact stateful CPU/GPU tokenization for append-heavy agent sessions, motivated by measurements showing repeated long-prefix tokenization can become a front-end bottleneck as prompt-cache hit rates rise.
\end{itemize}
{\bfseries 読む際の境界}
\begin{itemize}[leftmargin=1.5em,itemsep=0.35em]
\item \textbf{分析上の留意点}: The systems gains are author-reported and depend on analytical/microbenchmark assumptions, workload traces, serving stack, and hardware.
\end{itemize}
{\bfseries Primary source}
\begin{itemize}[leftmargin=1.5em,itemsep=0.25em]
\item \url{http://arxiv.org/abs/2607.10987v1}
\item \url{http://arxiv.org/abs/2607.23933v2}
\item \url{http://arxiv.org/abs/2607.29678v3}
\end{itemize}
{\scriptsize\color{SurveyMuted}Source-bound record: \texttt{evidence:SP-2026-M07:series-agent-serving-july-aa81e615c6}.}
\end{technicalnote}
\medskip
